\documentclass[aps,prl,twocolumn,superscriptaddress,showpacs,floatfix]{revtex4-2}

% =============================================================================
% PACKAGES
% =============================================================================
\usepackage{amsmath,amssymb,amsfonts}
\usepackage{graphicx}
\usepackage[utf8]{inputenc}
\usepackage[T1]{fontenc}
\usepackage{xcolor}
\usepackage{hyperref}
\usepackage{booktabs}
\usepackage{siunitx}
\usepackage{bm}

% =============================================================================
% DOCUMENT SETTINGS
% =============================================================================
\hypersetup{
    colorlinks=true,
    linkcolor=blue,
    citecolor=blue,
    urlcolor=blue
}

% =============================================================================
% CUSTOM COMMANDS
% =============================================================================
\newcommand{\Tc}{T_c}
\newcommand{\lambdaeph}{\lambda_{\text{ep}}}
\newcommand{\omegalog}{\omega_{\text{log}}}

% =============================================================================
% TITLE AND AUTHORS
% =============================================================================
\title{On the Necessity of Room-Temperature Superconductivity in Engineered Strong-Coupling Systems}

% [AUTHOR LIST - TO BE COMPLETED]
% \author{Author One}
% \affiliation{Affiliation One}
% \email{corresponding.author@institution.edu}
% 
% \author{Author Two}
% \affiliation{Affiliation Two}
% 
% \author{Author Three}
% \affiliation{Affiliation Three}

\date{\today}

% =============================================================================
% ABSTRACT
% =============================================================================
\begin{abstract}
Superconductivity at ambient conditions has long been considered the ``holy grail'' of condensed matter physics. 
Here we present a formal proof of the \textit{necessity} of room-temperature superconductivity (RTSC) in certain classes of engineered materials with strong electron-phonon coupling and reduced dimensionality. 
By combining theoretical analysis, computational predictions, and experimental validation, we demonstrate that...

[\textbf{Note:} Full abstract to be completed upon integration of A/B/C subagents]

\end{abstract}

% [PACS CODES - TO BE DETERMINED]
% \pacs{74.20.-z, 74.62.-c, 74.70.-b}

\maketitle

% =============================================================================
% INTRODUCTION
% =============================================================================
\section{Introduction}

The discovery of superconductivity in mercury by Kamerlingh Onnes in 1911~\cite{onnes1911} initiated a century-long quest for materials exhibiting zero electrical resistance at ever-higher temperatures. 
The BCS theory~\cite{bardeen1957}, developed in 1957, established the microscopic foundation for conventional superconductivity through electron-phonon coupling, predicting an isotope effect and an upper critical temperature limit ($\Tc \sim 30$--$40$~K) for simple metals---the so-called ``McMillan limit''~\cite{mcmillan1968}.

However, subsequent discoveries dramatically revised this picture. 
High-$\Tc$ cuprate superconductors, discovered by Bednorz and M\"{u}ller in 1986~\cite{bednorz1986}, achieved transition temperatures exceeding $90$~K in YBa$_2$Cu$_3$O$_{7-\delta}$ (YBCO), far above the McMillan limit. 
More recently, hydrides under high pressure have shattered records, with H$_3$S reaching $203$~K at $155$~GPa~\cite{drozdov2015} and LaH$_{10}$ achieving $250$--$280$~K at $170$--$200$~GPa~\cite{somayazulu2019,drozdov2019}.

While most research has focused on the \textit{sufficiency} of specific mechanisms for achieving high-$\Tc$ superconductivity, the question of \textit{necessity} has received less attention: \textbf{Under what conditions must room-temperature superconductivity exist?}
This framing shifts the perspective from serendipitous discovery to systematic design.

% =============================================================================
% THEORETICAL FRAMEWORK
% =============================================================================
\section{Theoretical Framework}

\subsection{Preliminaries and Definitions}

\begin{definition}[Strong Electron-Phonon Coupling Regime]
A material is in the strong electron-phonon coupling regime when the dimensionless coupling parameter $\lambda$ satisfies $\lambda > \lambda_c$, where $\lambda_c$ is a critical threshold determined by...
\end{definition}

[\textbf{To be imported from A路}]

\subsection{The Necessity Theorem}

\begin{theorem}[Necessity of RTSC]
Under conditions (i)--(iii), room-temperature superconductivity ($\Tc > 300$~K) is a necessary consequence.
\end{theorem}

[\textbf{Proof and detailed conditions to be imported from A路}]

% =============================================================================
% MATERIAL PREDICTIONS
% =============================================================================
\section{Material Predictions}

[\textbf{To be imported from B路}]

% =============================================================================
% EXPERIMENTAL VALIDATION
% =============================================================================
\section{Experimental Validation}

[\textbf{To be imported from C路}]

% =============================================================================
% DISCUSSION
% =============================================================================
\section{Discussion}

[\textbf{To be completed upon integration}]

% =============================================================================
% CONCLUSION
% =============================================================================
\section{Conclusion}

[\textbf{To be completed upon integration}]

% =============================================================================
% ACKNOWLEDGMENTS
% =============================================================================
\begin{acknowledgments}
[To be completed]
\end{acknowledgments}

% =============================================================================
% REFERENCES
% =============================================================================
\bibliography{references}

% =============================================================================
% APPENDIX (SUPPLEMENTARY DERIVATIONS)
% =============================================================================
\appendix

\section{Detailed Mathematical Derivations}

[\textbf{To be imported from A路}]

\section{Computational Methodology}

[\textbf{To be imported from B路}]

\section{Experimental Data Tables}

[\textbf{To be imported from C路}]

\end{document}
